% Canonical templates and operations over finite-dimensional coefficient algebras.
\subsection{C.3: finite local coefficient algebras over an algebraically closed field}

Let $k$ be algebraically closed of characteristic $p>2$, and let $R$ be a
nonzero finite-dimensional commutative local $k$-algebra. Fix a $k$-basis
$a_1=1,a_2,\ldots,a_c$ and its multiplication constants, and write
$\rho:R\to k$ for its residue map. The data are effective in the field-access
model of the preceding subsection. If the residue map is not supplied, it
is computed as follows: multiplication by $a_i$ on $R$ has the single
eigenvalue $\rho(a_i)$, because its difference from that scalar is
nilpotent. Factorization of its characteristic polynomial recovers this
unique eigenvalue, without dividing a trace by $c$.

Let $G_R=O_{m,k}\times_k R$ for the split orthogonal group. All representations
below are rational $G_R$-representations whose underlying $R$-modules are
finite free. Their duality is
\[
                 M^*=\operatorname{Hom}_R(M,R).
\]
In particular it is not the $k$-linear dual. No Frobenius or Gorenstein
hypothesis on $R$ is made. Put $H=k[O_{m,k}]$ and use the finite coalgebras
$C_d\subset H$ of Theorem~\ref{C:Om-canonical-templates}. Coefficient degree
over $R$ means degree in the filtration $R\otimes_k C_d$.

\begin{theorem}[Proved: canonical templates over every finite local coefficient algebra]
\label{C:artinian_algclosed_templates}
For every $n\geq1$ and $d\geq0$, a finite algorithm constructs polynomial
matrix templates over $k$, with polynomial/eigenvalue-order domains, for
exactly one representative of each indecomposable rational $G_R$-module
of $R$-rank $n$ and minimal coefficient degree $d$.
The templates encode the $R$-action and the rational group action on the
underlying $cn$-dimensional $k$-space. The algorithm recovers an $R$-basis
of every specialized member and its coaction matrices over $R$.

The union of these templates gives a parameter set $\mathscr I_R$ with no
remaining isomorphism quotient. $R$-duality is computable on this set.
The symmetric and alternating types are specified by finite Boolean
polynomial conditions on each canonical domain; every self-dual
indecomposable has exactly one of those types.
\end{theorem}
\begin{proof}
Write $a_i a_j=\sum_l\mu_{ij}^l a_l$. Let $Q_i$ be multiplication by
$a_i$ on the fixed basis of $R$, and set
\[
                       Q_i^{(n)}=I_n\otimes Q_i.
\]
Put $N=cn$, and let $u_1,\ldots,u_r$ be the fixed basis of $C_d^*$ with
multiplication and unit constants $\gamma_c^{ab},e_c$ as in (C.3.T2).
An $R\otimes_k C_d^*$-module on $k^N$ is a tuple
\[
       (B_1,\ldots,B_c,T_1,\ldots,T_r)
\]
subject to the polynomial equations
\[
 \begin{gathered}
 B_1=I_N,\qquad B_iB_j=\sum_l\mu_{ij}^l B_l,
       \qquad B_iT_a=T_aB_i,\\
 T_aT_b=\sum_l\gamma_l^{ab}T_l,
       \qquad \sum_a e_aT_a=I_N .
 \end{gathered}                                           \tag{C.3.A1}
\]
Its underlying $R$-module is free of rank $n$ exactly when
\[
 \exists S\in M_N(k)\quad
     B_iS=SQ_i^{(n)}\ (1\leq i\leq c),\qquad \det S\ne0.
                                                               \tag{C.3.A2}
\]
This is an existence statement about a basis, not a new orbit quotient.
Lemma~\ref{C:canonical-polynomial-elimination} eliminates $S$ and gives
an explicit Boolean polynomial condition on the entries of the tuple.

Apply Lemma~\ref{C:canonical-tuples} to the $c+r$ matrices of size $N$,
and impose (C.3.A1) and the eliminated condition (C.3.A2) on each domain.
A $k$-linear isomorphism intertwining the full tuple is precisely an
$R$-linear group-equivariant isomorphism. Thus the canonical list classifies
exactly the required objects. The equivalence with rational coactions is
also direct: for an $R$-module $M$,
\[
 M\otimes_R(R\otimes_k H)=M\otimes_k H,
\]
and $R$-linearity of the right side's coaction says exactly that the
$T_a$ commute with the $B_i$.

Require the coefficients outside $C_{d-1}$ not all to vanish when $d>0$.
This is the same minimal-degree test as over the field: the corresponding
$T_a$ are zero or nonzero independently of the $k$-basis. Since the
filtration is nested and exhausts $H$, it excludes all duplicates between
different $d$. Different $R$-ranks have different $k$-dimensions.

The joint centralizer of all $B_i,T_a$ is the endomorphism algebra in this
category. Any idempotent splits the object as an $R$-linear direct sum.
The two underlying modules are direct summands of $R^n$, hence finite
projective and therefore free over the local ring $R$. Accordingly the
nontrivial-idempotent test (C.3.T5), applied to the full tuple, is exactly
the decomposability test in the present category. Polynomial elimination
and complementation give its indecomposable locus. The canonical-record
and parameter-slot conventions of Lemma~\ref{C:canonical-tuples} remain
unchanged, so the resulting indices have no duplicates.

For a specialized member, equations $B_iS=SQ_i^{(n)}$ are linear in $S$.
Compute a basis $S_1,\ldots,S_t$ of their solution space. The free-module
condition means that $\det(\sum_j z_jS_j)$ is a nonzero polynomial.
A grid with $N+1$ distinct elements of $k$ in each variable contains a
point where it is nonzero. This gives an invertible $S$ by finitely many
tests. After conjugation by $S$, the $R$-action is the fixed regular
one. Each block of $S^{-1}T_aS$ is multiplication by a unique element of
$R$; its value on $1\in R$ recovers that element. Thus one obtains
matrices $\tau_a\in M_n(R)$ and the coaction
$\sum_a\tau_a\otimes b_a$ in an actual $R$-basis.

Use the antipode constants in (C.3.T2) to form the $R$-dual matrices
\[
                    \tau_a^*=\sum_b s_{ab}\tau_b^{\mathsf t}.
\]
Encode them again over $k$, together with the regular $R$-action, and
canonicalize the full tuple. This computes the unique dual index.
The antipode preserves $C_d$ and squares to the identity, so duality
preserves minimal coefficient degree.

For the polynomial form-type conditions one may keep $S$ as an auxiliary
variable. Introduce the entries of $\tau_a$ as $c$ scalar coordinates in
each position, and impose
\[
 T_aS=S\,\operatorname{Reg}(\tau_a),\qquad \det S\ne0,
 \qquad B_iS=SQ_i^{(n)},
\]
where $\operatorname{Reg}(\tau_a)$ replaces an entry of $R$ by its
multiplication matrix. Introduce $B\in M_n(R)$ and impose
\[
 B\tau_a=\tau_a^*B\ (\forall a),\qquad
 B^{\mathsf t}=B\ \hbox{or}\ B^{\mathsf t}=-B,
             \qquad \det(\rho(B))\ne0.                  \tag{C.3.A3}
\]
All are polynomial conditions over $k$ in the scalar coordinates.
The last inequality is equivalent to $\det_R B\in R^\times$, hence
exactly to nondegeneracy. Eliminate the auxiliary variables by
Lemma~\ref{C:canonical-polynomial-elimination}. This gives the asserted
polynomial type loci without making any symbolic choice of an $R$-basis.
For an actual form, extract a basis as above, solve the linear conditions
on $B$, and apply the determinant-grid test to its residue matrix.

Finally the category is additive and idempotent complete. Each endomorphism
algebra is finite-dimensional over $k$, so every object has a finite
Krull--Schmidt decomposition and each indecomposable endomorphism algebra
is local with nilpotent radical. Its residue division algebra is $k$:
every element is algebraic over $k$, and a split polynomial cannot
annihilate a non-scalar element of a division algebra. The $R$-dual
adjoint involution fixes $k$, hence induces the identity on this residue
field. Quebbemann--Scharlau--Schulte
\cite[Propositions 2.4--2.5 and Theorem 3.3, pp.271--273]{C-QSS1979}
therefore give exactly one of the symmetric and alternating types on each
self-dual indecomposable. Their hypotheses have been checked for this
$R$-dual category, and do not require an abelian category or a Gorenstein
coefficient algebra. This proves every assertion.
\end{proof}

\begin{theorem}[Proved: the full split orthogonal ring over finite local coefficient algebras]
\label{C:artinian_algclosed_WG}
Partition $\mathscr I_R$ into symmetric self-dual indices
$\mathscr I_{R,\mathrm s}$, alternating self-dual indices
$\mathscr I_{R,\mathrm a}$, and non-self-dual pairs. Choose the smaller
index of each pair in the fixed canonical order, giving
$\mathscr I_{R,\mathrm h}$. Then
$WG_{\mathrm{sp}}(B_{G_R})$ is free abelian on
\[
 (U_i,q_i)\quad(i\in\mathscr I_{R,\mathrm s}),\qquad
 H(U_i)\quad(i\in\mathscr I_{R,\mathrm a}\cup
                         \mathscr I_{R,\mathrm h}).
\]
Forgetting forms is injective into the split Green ring of finite free
$R$-representations, and sends these generators to
\[
 [U_i],\qquad 2[U_i],\qquad [U_i]+[U_{i^*}],
\]
respectively. All multiplication coefficients and exterior powers of
finite virtual inputs are computable in finitely many exact operations.
Tensor degrees are bounded by $d+e$, and the $j$th exterior-power degree
by $jd$.
\end{theorem}
\begin{proof}
The preceding proof verifies the three hypotheses of \S3 of
\cite{C-QSS1979}: idempotent completeness, Krull--Schmidt decompositions
with local endomorphism rings, and completeness for the endomorphism
radicals. The radicals are nilpotent here. Reduction and transfer,
\cite[Theorem 2.2, Corollary 2.3 and Theorems 3.2--3.3,
pp.269--273]{C-QSS1979}, classify forms on each self-dual isotypic part
by multiplicity forms over $k$. Symmetric multiplicity forms over an
algebraically closed field of characteristic different from two are
classified by dimension; alternating multiplicity forms are classified
by even dimension. Non-self-dual pairs are hyperbolic. This gives the
three stated generators, with disjoint supports in the indecomposable
basis, and proves injectivity of forgetting forms. These are the split
orthogonal-sum Grothendieck groups of the source's p.267, not its later
exact-sequence quotients.

Extract $R$-bases and coaction matrices as in the first theorem. Tensor
products are now tensor products over $R$, and exterior powers are exterior
powers over $R$. Their coefficient matrices are products and minors with
entries in $R\otimes_k H$. Use the multiplication table of $R$ and the
fixed normal forms of $H$ to express them in $R\otimes C_{d+e}$ or
$R\otimes C_{jd}$. No division by $j!$ is used. Encode their full
$k$-tuples and compute their indecomposable direct sums. Each summand is
again free over $R$; recanonicalize it at its least coefficient degree.
The resulting multiplicities recover the orthogonal coefficients by the
same rule (C.3.T9): divide alternating-type multiplicities by two and
retain one coefficient from each equal dual pair. Injective forgetting
makes the rule exact. Extend products by bilinearity and exterior powers
by $\lambda_t(-x)=\lambda_t(x)^{-1}$ coefficientwise. Each requested
coefficient uses only finitely many computations.
\end{proof}

\begin{corollary}[Proved: all finite-dimensional coefficient algebras over $k$]
\label{C:finite_algclosed_algebra_WG}
For every finite-dimensional commutative $k$-algebra $R$, with $k$
algebraically closed of characteristic $p>2$, the split ring
$WG_{\mathrm{sp}}(B_{O_{m,R}})$ has the same explicit finite-template
classification and finite operation procedures, applied separately to
the finitely many local factors of $R$. No reducedness hypothesis is
required.
\end{corollary}
\begin{proof}
Write the Artin decomposition $R=\prod_{j=1}^sR_j$. It is computable by
applying the canonical direct-sum algorithm to the commuting multiplication
matrices of the regular $R$-module: their indecomposable projective summands
are the primitive ideals $e_jR$, with their units recovered as the values
of the corresponding projections on $1$. Commutativity makes these
orthogonal idempotents central and gives the local factors $R_j$.
Finite projective representations and their $R$-duals decompose along
these idempotents. Symmetric forms and isometries decompose as well.
Hence the required ring is the finite product of the rings over $R_j$,
with componentwise tensor and exterior operations. The two theorems apply
to every factor.

Over each $R_j$, every nonsingular symmetric form on a free module of
rank $m$ is isometric to the standard one. First diagonalize its residue
form over $k$ and lift that change of basis. The resulting Gram matrix is
$I+N$ with entries of $N$ in the nilpotent maximal ideal. A finite polynomial
inverse square root, obtained recursively using only division by two,
gives $g$ with $g^{\mathsf t}(I+N)g=I$. Thus the assertion also covers
orthogonal groups defined by any specified nonsingular rank-$m$ symmetric
form over $R$, after this explicit change of coordinates.
\end{proof}

\begin{proposition}[Proved: a non-Gorenstein coefficient example]
\label{C:artinian_algclosed_example}
Let $R=k[a,b]/(a,b)^2$, where $k$ is algebraically closed of
characteristic $p>2$. The preceding construction applies
to $R$, although the ordinary $k$-dual of its regular module is not
free over $R$. In the basis $(1,a,b)$ its regular multiplication matrices
are
\[
 Q_a=\begin{pmatrix}0&0&0\\1&0&0\\0&0&0\end{pmatrix},\qquad
 Q_b=\begin{pmatrix}0&0&0\\0&0&0\\1&0&0\end{pmatrix}.
                                                               \tag{C.3.A4}
\]
On $R^2$ put
\[
 P=\begin{pmatrix}1&a\\b&1\end{pmatrix},\quad
 D=\operatorname{diag}(1,-1),\quad
 T=PDP^{-1},\quad B=P^{-\mathsf t}P^{-1}.
                                                               \tag{C.3.A5}
\]
Then $T^2=I$, $T^{\mathsf t}BT=B$, and $\det_R B=1$.
Thus $T$ defines a rational $O_1=\mu_2$ representation with an
invariant nonsingular symmetric $R$-form. Its $R$-dual action is
$T^{\mathsf t}$, with the regular $R$-action retained; its exterior
square is the sign representation.
\end{proposition}
\begin{proof}
The products $Q_a^2,Q_aQ_b,Q_bQ_a,Q_b^2$ vanish, so (C.3.A4)
faithfully encodes the multiplication table. A linear map on $R$
commutes with both matrices exactly when it is multiplication by its
value on $1$. In rank $n$, each of its $n^2$ blocks has this form,
giving a centralizer of dimension $3n^2$ and the stated entry-recovery
procedure.

The joint kernel of $Q_a,Q_b$ has dimension two. The joint kernel of
$Q_a^{\mathsf t},Q_b^{\mathsf t}$ has dimension one. Consequently
$\operatorname{Hom}_k(R,k)$ is not isomorphic to $R$; since it has
the same $k$-dimension as $R$, it is not a free $R$-module. This also
shows that $R$ is not Gorenstein: its $k$-dual is its dualizing module,
and such an Artin local algebra is Gorenstein precisely when that
module is free of rank one
\cite[Tags 0AX1 and 0DW6, Lemmas 47.16.1 and 47.21.4]{C-StacksArtinianDual}.
In contrast, evaluation at $1$ identifies
$\operatorname{Hom}_R(R,R)$ with the regular module $R$.

Writing $P=I+N$ gives $N^2=0$, hence $P^{-1}=I-N$.
The three matrix identities in the statement follow at once from
$D^2=I=D^{\mathsf t}D$ and $\det P=1$.
More explicitly,
\[
 T=\begin{pmatrix}1&-2a\\2b&-1\end{pmatrix},\qquad
 B=\begin{pmatrix}1&-a-b\\-a-b&1\end{pmatrix}.
\]
The two coaction coefficients are $(I+T)/2$ and $(I-T)/2$;
they are complementary idempotents, which verifies the rational
$\mu_2$ coaction. Its coefficient degree is one. The antipode
fixes the coordinate of $\mu_2$, so its $R$-dual has action
$T^{\mathsf t}$. Finally $\det_R T=-1$, proving the exterior-square
assertion.

The bounded executable certificate
\texttt{code/C\_finish\_artinian\_algclosed.py} performs these checks
over $\mathbb F_p$ for $p=3,5,7,11$; the identities therefore also
give examples after passage to their algebraic closures. It checks
all $36$ regular-basis products, the centralizers for free ranks
$1,2,3$ in each characteristic, and recovery of all $168$ resulting
centralizer basis elements. In each characteristic it then scrambles
the rank-two example by an invertible $k$-matrix, solves the regular
action's linear intertwiner equations to recover an $R$-basis, and
checks the recovered coaction, its $R$-dual, invariant form, residue
determinant, and exterior square. It separately checks the distinct
socle dimensions of the regular modules and their ordinary $k$-duals.
The log records these $12$ centralizer cases and $4$ basis-recovery
cases. This verifies the bounded example, not the unbounded template
algorithm.
\end{proof}

\paragraph{Scope.}
The coefficient algebras are finite-dimensional over a specified
algebraically closed field. The result does not assert an algorithm for
arbitrary nonnoetherian bases or for local rings of positive dimension.
It adds a nonfield case to the field construction while preserving the
explicit field/factorization/order access assumptions and the finite
per-rank, per-coefficient-degree scope.
